%%=============================================================================
%% CAPITULO 2 -- REFERENCIAL TEORICO
%%
%% Esta e a parte do trabalho que mais demanda dominio das regras de citacao
%% (diretas e indiretas). Reveja a NBR 10520 antes de escrever.
%%=============================================================================

\section{Referencial teórico}

O método de condução desta revisão está descrito na Seção de procedimentos
metodológicos. Este capítulo revisa a literatura que fundamenta cada decisão
arquitetural do sistema avaliado neste trabalho, seguindo a ordem do próprio
pipeline de processamento --- da origem do paradigma de RAG à avaliação dos
resultados. As seções seguintes situam a recuperação híbrida na evolução do
RAG, posicionam a arquitetura de GraphRAG incremental adotada frente às
variantes já publicadas e fecham o capítulo com a metodologia de avaliação e
o posicionamento deste trabalho frente a pesquisas comparáveis.

\subsection{RAG: origem e evolução do paradigma}

O paradigma de Retrieval-Augmented Generation (RAG) combina um modelo de
linguagem com um mecanismo de recuperação sobre uma base documental externa
ao modelo --- não aos parâmetros aprendidos durante o treinamento, mas a uma
base consultada em tempo de inferência ---, condicionando a geração da
resposta ao conteúdo recuperado em vez de depender só do conhecimento
implícito nos parâmetros \parencite{lewis2020}. A formulação original acopla
um retriever denso (embeddings) a um gerador treinado conjuntamente, aplicado
a tarefas intensivas em conhecimento como perguntas e respostas em domínio
aberto.

%% --- Figura (minha, não a de exemplo do template) ---------------------------
\begin{figure}[htb]
  \caption{Esquema ilustrativo do pipeline de recuperação e geração avaliado}
  \label{fig:pipeline-rag}
  \centering
  \includegraphics[width=0.6\textwidth]{figuras/exemplo-figura.png}
  \fonte{Elaboração própria (\the\year).}
\end{figure}

A Figura~\ref{fig:pipeline-rag} ilustra, de forma simplificada, o percurso que
uma pergunta percorre no sistema avaliado: consulta, recuperação (densa,
esparsa ou por grafo), geração da resposta a partir da evidência recuperada
e, na condição verificada, uma etapa adicional de auditoria antes da resposta
final. É apenas um esquema de trabalho desta fase do projeto; o diagrama
definitivo será redesenhado quando a arquitetura estiver congelada.

\textbf{Citação direta com mais de 3 linhas} --- recuo de 4~cm, fonte 10,
espaçamento simples e sem aspas, no idioma original do texto fonte:

Assim, segundo Lewis et al., sintetizando a motivação original do RAG:

\begin{citacao}
Large pre-trained language models have been shown to store factual knowledge
in their parameters, and achieve state-of-the-art results when fine-tuned on
downstream NLP tasks. However, their ability to access and precisely
manipulate knowledge is still limited, and hence on knowledge-intensive
tasks, their performance lags behind task-specific architectures.
\parencite{lewis2020}
\end{citacao}

\textbf{Citação indireta} --- texto baseado na obra consultada, com a ideia
transcrita em palavras próprias:

Para \citeonline{edge2024}, a limitação central do RAG convencional não está
na recuperação de um fato isolado, e sim em perguntas ``globais'', que pedem
uma síntese sobre toda a coleção de documentos; a solução proposta constrói
um grafo de conhecimento a partir dos documentos-fonte e pré-gera resumos de
comunidades de entidades relacionadas, combinando-os em tempo de consulta
para responder a esse tipo de pergunta --- com ganhos relatados de
abrangência e diversidade da resposta frente a um RAG convencional.

%% --- Tabela (minha, não a de exemplo do template) ---------------------------
\begin{table}[htb]
  \caption{Condições de recuperação comparadas neste trabalho}
  \label{tab:condicoes-recuperacao}
  \centering
  \begin{tabular}{lccc}
    \toprule
    \textbf{Condição} & \textbf{Índice vetorial} & \textbf{Índice esparso (BM25)} & \textbf{Grafo de conhecimento} \\
    \midrule
    RAG denso          & Sim & Não & Não \\
    RAG denso-esparso   & Sim & Sim & Não \\
    GraphRAG incremental & Sim & Não & Sim \\
    \bottomrule
  \end{tabular}
  \fonte{Elaboração própria (\the\year).}
\end{table}

A Tabela~\ref{tab:condicoes-recuperacao} resume as três condições de
recuperação comparadas na fase confirmatória deste trabalho, antes da
ampliação exploratória para um número maior de condições.
